







\documentclass[11pt,a4paper]{article}
\usepackage{luatexja}
\usepackage[margin=1in]{geometry}
\usepackage{amsmath,amssymb}
\usepackage{microtype}
\usepackage{booktabs}
\usepackage{array}
\usepackage{xcolor}
\usepackage{pgfplots}
\pgfplotsset{compat=1.17}
\usepgfplotslibrary{fillbetween}
\usepackage[numbers,sort&compress]{natbib}
\usepackage{fancyhdr}
\definecolor{medusablue}{HTML}{2A78D6}
\definecolor{c1}{HTML}{2A78D6}
\definecolor{c2}{HTML}{EB6834}
\definecolor{c3}{HTML}{1BAF7A}
\definecolor{muted}{HTML}{52514E}
\usepackage[colorlinks=true,linkcolor=medusablue,citecolor=medusablue,urlcolor=medusablue]{hyperref}

\pagestyle{fancy}
\fancyhf{}
\renewcommand{\headrulewidth}{0.4pt}
\lhead{\small\scshape A Preprint \textemdash{} Medusa Lab 2026-W39}
\rhead{\small\scshape 2026-09-24}
\cfoot{\thepage}
\setlength{\headheight}{14pt}

\title{Naming Game model of lexical convention formation: an in-silico study}
\author{%
  Medusa Lab agents: Scout, Analyst, Coder, Writer and Reviewer\\
  {\small Autonomous multi-agent research lab \quad Director: 所長}}
\date{2026-09-24\\[2pt] {\small Weekly theme: LLMはオノマトペを接地・創発できるか}}

\begin{document}
\maketitle
\thispagestyle{fancy}

\begin{abstract}
Medusa Lab's weekly research theme, “LLMはオノマトペを接地・創発できるか”, was mapped onto the Naming Game model of lexical convention formation. Agents negotiate names for an object through pairwise interactions until a shared convention emerges. We simulated the model while sweeping the number of agents $N$ over 32, 64, 128, 256, with 5 independent replicates per setting and 95\% confidence intervals. On the complete graph, convergence time grows as $N^{1.39}$ ($R^{2}$=0.97); the mean-field prediction is $N^{1.5}$. On the small-world network the fitted exponent is 2.09, and the peak memory grows as $N^{1.01}$ versus $N^{1.40}$ on the complete graph. Of the 2 in-silico hypotheses, 2 are supported by the simulation. The scaling exponent quantifies how costly it is for larger communities to agree on a convention. The study was run end-to-end by autonomous agents in offline mode; its results are hypotheses for further work, and companion proposals describe experiments that need human participants.
\end{abstract}

\noindent\textbf{Keywords:} computational linguistics, LLM, sound symbolism, language evolution, agent-based simulation

\medskip
\noindent\fcolorbox{muted}{white}{\parbox{\dimexpr\linewidth-2\fboxsep-2\fboxrule\relax}{\small
\textbf{AI-generated manuscript.} This preprint was produced autonomously by the Medusa Lab multi-agent
system (Offline template mode (no language model)) and has not been peer-reviewed by humans. Every number comes from the archived
simulation (seed 20260901); verify before citing.}}

\section{Introduction}
This paper is the in-silico part of Medusa Lab's weekly research cycle. The Director chose the theme “LLMはオノマトペを接地・創発できるか”; the Analyst agent translated it into a stylised, simulable mechanism, the Naming Game model of lexical convention formation (computational linguistics / language evolution).

Shared vocabularies can emerge without central coordination, as first shown by agent-based language games \citep{steels1995selforganizing}. The minimal Naming Game captures this self-organisation with a sharp transition towards consensus and a characteristic $N^{3/2}$ scaling of the convergence time in well-mixed populations \citep{baronchelli2006sharp}. Contact topology strongly affects the dynamics \citep{dallasta2006nonequilibrium}, and experiments with human groups confirm that conventions emerge spontaneously in networked populations \citep{centola2015spontaneous}.

The Scout agent identified the following open questions:
\begin{itemize}
  \item Here, to address these limitations, we present MultiMedQA, a benchmark combining six existing medical question answering datasets spanning professional medicine, research and consumer queries and a new dataset of medical questions searched online, HealthSearchQA.
  \item However, they can also be harnessed for various attacks (particularly user-level attacks) due to their human-like reasoning abilities.
  \item Joan Bybee and her colleagues present a new theory of the evolution of grammar that links structure and meaning in a way that directly challenges most contemporary versions of generative grammar.
\end{itemize}

Our contributions are:
\begin{itemize}
  \item a reproducible simulation of the model with a systematic parameter sweep;
  \item tests of 2 in-silico hypotheses;
  \item 2 human-experiment proposal(s) that confront the model with human behaviour.
\end{itemize}

\section{Related work}
Shared vocabularies can emerge without central coordination, as first shown by agent-based language games \citep{steels1995selforganizing}. The minimal Naming Game captures this self-organisation with a sharp transition towards consensus and a characteristic $N^{3/2}$ scaling of the convergence time in well-mixed populations \citep{baronchelli2006sharp}. Contact topology strongly affects the dynamics \citep{dallasta2006nonequilibrium}, and experiments with human groups confirm that conventions emerge spontaneously in networked populations \citep{centola2015spontaneous}.

Recent work retrieved by the Scout agent includes the following: \citet{imai2014sound} study The sound symbolism bootstrapping hypothesis for language acquisition and language evolution; \citet{singhal2023large} study Large language models encode clinical knowledge; \citet{hu2021lora} study LoRA Fine-Tuning of a 3B Code LLM for Algorithmic Efficiency; \citet{christian2025neuromodulatory} study Neuromodulatory Control Networks (NCNs): A Biologically Inspired Architecture for Dynamic LLM Processing; \citet{yao2024survey} study A survey on large language model (LLM) security and privacy: The Good, The Bad, and The Ugly.

\section{Model}
Each agent holds an inventory of candidate names. At every step a random speaker utters a name from its inventory (inventing a new one if empty) to a random neighbour. If the hearer already knows the name, the interaction succeeds and both agents discard all competing names; otherwise the hearer adds the name to its inventory.

Assumptions:
\begin{itemize}
  \item Agents have unbounded memory and no preference among names.
  \item Interactions are pairwise, random and sequential.
  \item A single object (meaning) is being named; there is no homonymy.
\end{itemize}

The governing equations are given in Equations (1)--(3) and the parameters in Table~\ref{tab:params}.

\begin{equation}\label{eq:1}
N_w(t) = \sum_{i=1}^{N} |\mathcal{I}_i(t)|
\end{equation}
\begin{equation}\label{eq:2}
t_{\mathrm{conv}} \sim N^{\alpha}, \quad \alpha_{\mathrm{MF}} = 3/2
\end{equation}
\begin{equation}\label{eq:3}
N_w^{\max} \sim N^{\beta}, \quad \beta_{\mathrm{MF}} = 3/2
\end{equation}
\begin{table}[ht]
\centering\small
\caption{Model parameters.}\label{tab:params}
\begin{tabular}{l l p{0.36\linewidth} l p{0.2\linewidth}}
\toprule
Symbol & Name & Description & Default & Sweep\\
\midrule
$N$ & population & number of agents & 256 & 32, 64, 128, 256\\
$k$ & sw\_degree & degree of the small-world ring lattice & 6 & --\\
$p$ & sw\_rewire & Watts-Strogatz rewiring probability & 0.1 & --\\
\bottomrule
\end{tabular}
\end{table}
\section{Methods}
We swept the number of agents $N$ over 32, 64, 128, 256. For every setting we ran 5 independent replicates with seeds derived from the base seed 20260901, and report means with 95\% t-intervals.

The simulation (coder/sim.py) was executed in Medusa Lab's sandbox (isolation: strict); the full run took 1.7 seconds. The code, parameters and raw data tables are archived with this paper, so every number and figure can be regenerated.

\paragraph{Hypotheses tested in silico.}
\begin{itemize}
  \item[\textbf{H1}] In a well-mixed population the time to lexical consensus grows super-linearly with population size, with an exponent close to the mean-field value 3/2.
  \item[\textbf{H2}] Local small-world contact structure lowers the peak lexical memory but slows the final agreement compared with a well-mixed population.
\end{itemize}
\section{Results}
Figure~\ref{fig:fig_convergence}, Figure~\ref{fig:fig_success} summarise the parameter sweep.

\begin{itemize}
  \item On the complete graph, convergence time grows as $N^{1.39}$ ($R^{2}$=0.97); the mean-field prediction is $N^{1.5}$.
  \item On the small-world network the fitted exponent is 2.09, and the peak memory grows as $N^{1.01}$ versus $N^{1.40}$ on the complete graph.
  \item At N=256, consensus took 30.3 interactions per agent on the complete graph and 412.4 on the small-world network.
\end{itemize}

Table~\ref{tab:metrics} lists the summary metrics.

\textbf{H1} is supported: fitted exponent 1.39 vs. the expected 1.3-1.6 (mean field: 1.5).

\textbf{H2} is supported: peak-memory exponent 1.014 (small-world) vs 1.396 (complete); time per agent 412.41 vs 30.27.

\begin{figure}[ht]
\centering
\begin{tikzpicture}
\begin{axis}[
  width=0.84\linewidth, height=0.52\linewidth,
  xlabel={Population size $N$}, ylabel={Convergence time per agent $t_{conv}/N$},
  xmode=log, log ticks with fixed point, xtick=data,
  ymode=log, log ticks with fixed point,
  grid=major, major grid style={line width=0.2pt, draw=black!12},
  axis line style={draw=black!45}, tick style={draw=black!45},
  tick label style={font=\small}, label style={font=\small},
  legend style={font=\small, draw=none, fill=none, at={(0.5,-0.24)}, anchor=north, legend columns=-1},
  every axis plot/.append style={line width=1.1pt}, mark size=1.7pt,
]
\addplot[draw=none, forget plot, name path=lo0] table[x=N, y=t_mf_lo] {data/convergence.dat};
\addplot[draw=none, forget plot, name path=hi0] table[x=N, y=t_mf_hi] {data/convergence.dat};
\addplot[c1, fill opacity=0.15, draw=none, forget plot] fill between[of=hi0 and lo0];
\addplot[c1, mark=*, mark options={solid}] table[x=N, y=t_mf] {data/convergence.dat};
\addlegendentry{complete graph}
\addplot[draw=none, forget plot, name path=lo1] table[x=N, y=t_sw_lo] {data/convergence.dat};
\addplot[draw=none, forget plot, name path=hi1] table[x=N, y=t_sw_hi] {data/convergence.dat};
\addplot[c2, fill opacity=0.15, draw=none, forget plot] fill between[of=hi1 and lo1];
\addplot[c2, mark=*, mark options={solid}] table[x=N, y=t_sw] {data/convergence.dat};
\addlegendentry{small-world}
\end{axis}
\end{tikzpicture}
\caption{Time to lexical consensus per agent versus population size (mean and 95\% CI over 5 runs).}\label{fig:fig_convergence}
\end{figure}
\begin{figure}[ht]
\centering
\begin{tikzpicture}
\begin{axis}[
  width=0.84\linewidth, height=0.52\linewidth,
  xlabel={Interactions per agent}, ylabel={Communicative success rate},
  grid=major, major grid style={line width=0.2pt, draw=black!12},
  axis line style={draw=black!45}, tick style={draw=black!45},
  tick label style={font=\small}, label style={font=\small},
  legend style={font=\small, draw=none, fill=none, at={(0.5,-0.24)}, anchor=north, legend columns=-1},
  every axis plot/.append style={line width=1.1pt}, mark size=1.7pt,
]
\addplot[c1, mark=none] table[x=t, y=success_mf] {data/success.dat};
\addlegendentry{complete graph}
\addplot[c2, mark=none] table[x=t, y=success_sw] {data/success.dat};
\addlegendentry{small-world}
\end{axis}
\end{tikzpicture}
\caption{Success rate of interactions over time for $N=256$ (sliding window).}\label{fig:fig_success}
\end{figure}
\begin{table}[ht]
\centering\small
\caption{Summary metrics reported by the simulation.}\label{tab:metrics}
\begin{tabular}{p{0.3\linewidth} l p{0.46\linewidth}}
\toprule
Metric & Value & Description\\
\midrule
alpha\_\allowbreak{}complete & 1.386 & exponent of convergence time t\_conv $\sim$ $N^{\alpha}$ (complete graph)\\
alpha\_\allowbreak{}small\_\allowbreak{}world & 2.094 & exponent of t\_conv $\sim$ $N^{\alpha}$ (small-world network)\\
r2\_\allowbreak{}complete & 0.973 & $R^{2}$ of the log-log fit (complete graph)\\
peak\_\allowbreak{}exponent\_\allowbreak{}complete & 1.396 & exponent of peak memory $N_{\mathrm{w}}^{max}$ $\sim$ $N^{\beta}$\\
peak\_\allowbreak{}exponent\_\allowbreak{}small\_\allowbreak{}world & 1.014 & exponent of peak memory (small-world)\\
t\_\allowbreak{}per\_\allowbreak{}agent\_\allowbreak{}complete\_\allowbreak{}maxN & 30.27 & convergence time per agent at the largest N\\
t\_\allowbreak{}per\_\allowbreak{}agent\_\allowbreak{}small\_\allowbreak{}world\_\allowbreak{}maxN & 412.4 & convergence time per agent at the largest N (SW)\\
consensus\_\allowbreak{}rate & 1 & fraction of runs that reached consensus\\
\bottomrule
\end{tabular}
\end{table}
\section{Discussion}
\begin{itemize}
  \item The scaling exponent quantifies how costly it is for larger communities to agree on a convention.
  \item Sparse local contacts trade memory load against speed: agents juggle fewer names but competing local conventions survive longer.
  \item Online platforms that reshuffle contacts behave more like well-mixed populations, which predicts faster adoption of neologisms.
\end{itemize}

Several questions cannot be settled in silico. We therefore propose human experiments: E1 “Online naming-game experiment: does network structure speed up agreement on new words?”; E2 “Intuition vs simulation: can people foresee macro outcomes from micro rules?”. Their designs use this simulation to set conditions and expected effect sizes.

\subsection*{Proposed human experiments}
Questions that simulation alone cannot settle were written up for the Director as separate experiment protocols:
\begin{itemize}
  \item \textbf{E1}: Online naming-game experiment: does network structure speed up agreement on new words? ($n = 12$ per condition from a power analysis).
  \item \textbf{E2}: Intuition vs simulation: can people foresee macro outcomes from micro rules? ($n = 64$ per condition from a power analysis).
\end{itemize}
\section{Limitations}
\begin{itemize}
  \item The minimal Naming Game ignores meaning, memory decay and prestige-biased imitation.
  \item Population sizes are small; asymptotic exponents may differ from the fitted finite-size values.
  \item A single topology per class (complete graph, one small-world ensemble) was studied.
\end{itemize}

The paper was written by Medusa Lab's template writer (offline mode, no language model); the mapping from the theme to the model is a deliberate simplification, and hypothesis verdicts follow pre-specified numerical criteria.

\section{Conclusion}
Within the Naming Game model of lexical convention formation, on the complete graph, convergence time grows as $N^{1.39}$ ($R^{2}$=0.97); the mean-field prediction is $N^{1.5}$. In short: consensus takes disproportionately longer as the group grows. The companion human-experiment proposals outline how to test the mechanism with people.

\section*{Reproducibility statement}
The simulation script, its parameters, the raw data tables and the run log are archived next to this paper in the
Medusa Lab repository (cycle 2026-W39). The run used seed 20260901, took 1.7\,s and was
executed in the strict sandbox (code source: recipe).

\begin{thebibliography}{99}
\bibitem[Steels(1995)]{steels1995selforganizing} Luc Steels (1995). A self-organizing spatial vocabulary. \emph{Artificial Life 2(3):319-332}.
\bibitem[Baronchelli et al.(2006)]{baronchelli2006sharp} Andrea Baronchelli, Maddalena Felici, Vittorio Loreto, Emanuele Caglioti, Luc Steels (2006). Sharp transition towards shared vocabularies in multi-agent systems. \emph{Journal of Statistical Mechanics: Theory and Experiment, P06014}. \url{https://doi.org/10.1088/1742-5468/2006/06/P06014}
\bibitem[Dall'Asta et al.(2006)]{dallasta2006nonequilibrium} Luca Dall'Asta, Andrea Baronchelli, Alain Barrat, Vittorio Loreto (2006). Nonequilibrium dynamics of language games on complex networks. \emph{Physical Review E 74, 036105}.
\bibitem[Centola and Baronchelli(2015)]{centola2015spontaneous} Damon Centola, Andrea Baronchelli (2015). The spontaneous emergence of conventions: An experimental study of cultural evolution. \emph{PNAS 112(7):1989-1994}. \url{https://doi.org/10.1073/pnas.1418838112}
\bibitem[Imai and Kita(2014)]{imai2014sound} Mutsumi Imai, Sotaro Kita (2014). The sound symbolism bootstrapping hypothesis for language acquisition and language evolution. \emph{Philosophical Transactions of the Royal Society B Biological Sciences}. \url{https://doi.org/10.1098/rstb.2013.0298}
\bibitem[Singhal et al.(2023)]{singhal2023large} Karan Singhal, Shekoofeh Azizi, Tao Tu, S. Sara Mahdavi, JASON KAI WEI LEE, Hyung Won Chung et al. (2023). Large language models encode clinical knowledge. \emph{Nature}. \url{https://doi.org/10.1038/s41586-023-06291-2}
\bibitem[Hu et al.(2021)]{hu2021lora} J. Edward Hu, Yelong Shen, Phillip Wallis, Zeyuan Allen-Zhu, Yuanzhi Li, Shean Wang et al. (2021). LoRA Fine-Tuning of a 3B Code LLM for Algorithmic Efficiency. \emph{arXiv (Cornell University)}. \url{https://doi.org/10.48550/arxiv.2106.09685}
\bibitem[Christian(2025)]{christian2025neuromodulatory} Morgan, Michael Christian (2025). Neuromodulatory Control Networks (NCNs): A Biologically Inspired Architecture for Dynamic LLM Processing. \emph{Zenodo (CERN European Organization for Nuclear Research)}. \url{https://doi.org/10.5281/zenodo.17575540}
\bibitem[Yao et al.(2024)]{yao2024survey} Yifan Yao, Jinhao Duan, Kaidi Xu, Yuanfang Cai, Zhibo Sun, Yue Zhang (2024). A survey on large language model (LLM) security and privacy: The Good, The Bad, and The Ugly. \emph{High-Confidence Computing}. \url{https://doi.org/10.1016/j.hcc.2024.100211}
\bibitem[Kantartzis et al.(2011)]{kantartzis2011japanese} Katerina Kantartzis, Mutsumi Imai, Sotaro Kita (2011). Japanese Sound-Symbolism Facilitates Word Learning in English-Speaking Children. \emph{Cognitive Science}. \url{https://doi.org/10.1111/j.1551-6709.2010.01169.x}
\bibitem[Paradis et al.(2004)]{paradis2004analyses} Emmanuel Paradis, Julien Claude, Korbinian Strimmer (2004). APE: Analyses of Phylogenetics and Evolution in R language. \emph{Bioinformatics}. \url{https://doi.org/10.1093/bioinformatics/btg412}
\bibitem[William(1995)]{william1995evolution} Bybee, Joan L. and Perkins, Revere and Pagliuca, William (1995). The evolution of grammar: tense, aspect, and modality in the languages of the world. \emph{Choice Reviews Online}. \url{https://doi.org/10.5860/choice.32-4321}
\bibitem[Mufwene(2001)]{mufwene2001ecology} Salikoko S. Mufwene (2001). The Ecology of Language Evolution. \emph{Cambridge University Press eBooks}. \url{https://doi.org/10.1017/cbo9780511612862}
\bibitem[Nadkarni et al.(2011)]{nadkarni2011natural} Prakash M. Nadkarni, Lucila Ohno-Machado, Wendy W. Chapman (2011). Natural language processing: an introduction. \emph{Journal of the American Medical Informatics Association}. \url{https://doi.org/10.1136/amiajnl-2011-000464}
\bibitem[Team(2000)]{team2000language} R Core Team (2000). R: A Language and Environment for Statistical Computing. \emph{}. \url{https://doi.org/10.32614/r.manuals}
\bibitem[Lewis et al.(2020)]{lewis2020affordancecompiled} Patrick Lewis, Ethan Perez, Aleksandara Piktus, Fabio Petroni, Vladimir Karpukhin, Naman Goyal et al. (2020). Affordance-Compiled Intelligence: Observable-Only Cognitive Impedance Matching for No-Meta LLM-Integrated Systems. \emph{arXiv (Cornell University)}. \url{https://doi.org/10.48550/arxiv.2005.11401}
\end{thebibliography}
\end{document}
